\section{Introduction}
Adaptive optimizer states can be expensive even when useful gradient directions occupy a much smaller subspace than the parameter matrix. Low-rank optimizers exploit this observation by maintaining adaptive statistics in projected coordinates \cite{zhao2024galore,hao2024flora,chen2024fira,zhu2024apollo,liang2024osd,robert2024ldadam,rajabi2025subtrackpp}. The scientific question is how to update that subspace as training changes, while accounting for the gradient outside it and the coordinate dependence of the adaptive states.

GrassRoute maintains one right routing subspace per selected matrix. It updates the subspace through a Grassmann tangent step with momentum and QR retraction, applies Adam to projected coordinates, and updates the orthogonal residual with normalized SGD. The parameter remains unconstrained and receives a full-dimensional update. Persistent optimizer state can be much smaller than AdamW's, although routing adds computation.

A rigorous account must distinguish several issues. Normalizing each covariance sample changes the population being estimated. A subspace is invariant to a basis change but diagonal Adam is not. Unbiased gradients do not imply unbiased adaptive directions. A nonzero residual branch also prevents treating its spectral tail as a universally irreducible optimization error. Our theory addresses these issues explicitly. The fixed-covariance landscape has a principal-angle PL bound; stochastic tracking adds noise, drift, and momentum-lag budgets. Parameter convergence retains an adaptive-direction defect. Generalization is analyzed through the full stored state and through an alternative PAC-Bayes perturbation bound. Detailed proofs occupy Appendix~\ref{app:theory}.

The empirical evidence identifies useful operating regimes for selective adaptive state. Archived Burgers/PINN trajectories show a favorable training-loss--state trade-off, while the independent Cora comparison improves on AdamW and corrected periodic routing with substantially less state. These observations concern specific objectives and resource budgets; they do not make GrassRoute the highest-accuracy optimizer in every comparison. Section~\ref{sec:new-public-tasks} adds six public-data tasks under a frozen protocol and explicitly distinguishes primary-metric gains from Pareto-efficient memory--quality trade-offs. The complete attempted evaluation, including unfavorable and unfinished cases, is retained in Appendices~\ref{app:full-results} and~\ref{app:new-task-full}.

We supplement that archive with a representative-model matrix designed to cover each task or dataset named in the proposal, using seed 42 and 1000 or 2000 complete epochs. It fixes historical rank overrides and GrassRoute's diagonal covariance transport, uses validation data for selection, and includes the official SubTrack++ implementation. After 117 of 162 planned formal runs, we stopped the remaining high-cost queue and documented that post-hoc scope change. An independent follow-up corrects the local periodic comparator's omitted second-moment transport. New and archived algorithm versions remain separate. Tables report only completed formal runs; the coverage audit records pending or failed comparisons instead of filling them with predictions. Additional actual-implementation experiments test tracking against instantaneous and full-covariance EMA spectral estimators, and expose the noise--lag--state trade-off across the entire declared drift grid.

Our contributions are a detailed geometric and conditional convergence analysis, full-state and PAC-Bayesian generalization analyses with explicit scope, an auditable experimental pipeline, and a task-by-task assessment of accuracy, memory, and compute. The angular PL factor is related to existing Rayleigh-quotient theory \cite{alimisis2024geodesic}; it is not presented as a uniquely established empirical predictor of optimizer rankings.
